\chapter{Aprendizaje Supervisado: Modelado Predictivo}

\section{Diseño Experimental y Ecosistema de Modelos}
El objetivo principal (\textit{core}) de la investigación residía en la construcción de un motor de inferencia capaz de aproximar la valoración económica del mercado inmobiliario. 
Para asegurar el rigor, la evaluación se rigió bajo un paradigma estricto de \textit{Hold-out Test} (partición 80/20) asistido por \textbf{Cross-Validation} estructurada (K-Folds). Adicionalmente, todo el protocolo incluyó bloqueos de semillas (\texttt{RANDOM\_STATE=42}) para imposibilitar la fuga de datos o discrepancias en la reproducibilidad.

Dada la heterogeneidad y los picos de no-linealidad expuestos en la fase de clustering, se entrenó masivamente un ecosistema integrado por 12 algoritmos distintos:
\begin{enumerate}
    \item \textbf{Modelos de Regresión Lineal y Penalizados:} Regresión Lineal OLS (Baseline), Regresión Ridge (penalización $L_2$) y Regresión Lasso (penalización $L_1$).
    \item \textbf{Basados en Árboles:} Árbol de Decisión (CART).
    \item \textbf{Support Vector Machines:} LinearSVR y SVR con kernel Radial (SVR\_RBF).
    \item \textbf{Ensembles (Bagging y Boosting):} RandomForestRegressor, AdaBoost, XGBoost \cite{chen2016xgboost} y Bagging sobre regresores SVR.
    \item \textbf{Meta-Modelos (Ensamblajes de Alto Nivel):} Voting Regressor y Stacking Regressor (ensamblajes heterogéneos anidados).
\end{enumerate}

\section{Optimización Paramétrica (GridSearchCV)}
La extracción de máxima capacidad teórica a los algoritmos se orquestó a través de un proceso de optimización exhaustiva, \texttt{GridSearchCV}. El grid de búsqueda iteró sobre los principales parámetros de control de la arquitectura de red algorítmica: profundidades de árbol (\texttt{max\_depth}), márgenes duros y blandos para el caso SVM (\texttt{C} y \texttt{gamma}), y coeficientes de penalización ($\alpha$). Se usó el coeficiente de determinación ($R^2$) como métrica unificadora de convergencia en este espacio cartesiano paramétrico.

\section{Evaluación, Diagnóstico y Selección Final}
Más allá de las métricas puramente aisladas como el Error Absoluto Medio (MAE) o el $R^2$, el protocolo de auditoría incorporó una métrica penalizante diseñada en este estudio: el \textbf{\textit{Overfitting Gap}}. Esta dimensión calculaba la sustracción analítica entre la varianza explicada en la partición de entrenamiento (Train) y la de evaluación (Hold-out). Cualquier modelo superando un salto del 0.05 era clasificado como inestable (memorización de ruido o varianza de sesgo extrema) y descartado categóricamente.

\begin{table}[H]
    \centering
    \caption{Resultados del Hold-out Test para modelos estables (\textit{Gap} $< 0.05$)}
    \begin{tabular}{lcccc}
        \toprule
        \textbf{Modelo} & \textbf{Hold-out $R^2$} & \textbf{MAE (\$)} & \textbf{RMSE (\$)} & \textbf{Overfitting Gap} \\
        \midrule
        \textbf{SVR\_RBF} & \textbf{0.5898} & \textbf{54.29} & \textbf{194.85} & \textbf{-0.0026} \\
        Voting\_Assembly & 0.5866 & 55.15 & 193.95 & 0.0414 \\
        XGBoost & 0.5863 & 55.40 & 193.10 & 0.0287 \\
        LinearRegression & 0.5471 & 57.32 & 196.29 & -0.0199 \\
        \bottomrule
    \end{tabular}
\end{table}

\textbf{Modelo Ganador:} El algoritmo \textbf{Support Vector Regressor con Kernel Radial (SVR\_RBF)} se posicionó inequívocamente como la solución definitiva. Su kernel no lineal ($RBF$) probó ser la herramienta matemática perfecta para capturar las interacciones geolocalizadas complejas dentro del espacio PCA. Alcanzó un $R^2$ aproximado del 59\% con un MAE competitivo de \$54.29 dólares y un \textit{Overfitting Gap} negativo, certificando matemáticamente su sublime e intachable capacidad de generalización algorítmica ante nueva casuística.
